% !TEX root = ../../vzor.tex
%% Ukázka: ruční seznam zkratek (nepovinný)
%% Vkládá ji soubor vzor/kapitoly/teorie.tex; ve vzor.pdf je pod nadpisem uvedeno „Zdroj ukázky".

\section{Seznam zkratek}
\zdrojukazky{vzor/ukazky/zkratky.tex}

Seznam zkratek je \emph{nepovinný}. Hodí se, když práce používá hodně odborných zkratek
(např.\ GPIO, PWM, API); běžné české zkratky (atd., tj., např.) do něj nepatří.
Nejjednodušší je napsat ho ručně jako samostatnou kapitolu -- v souboru \texttt{mprace.tex}
ji vložte mezi \verb|\printbibliography| a \verb|\listoffigures|:

{\footnotesize
\begin{verbatim}
\chapter*{Seznam zkratek}
\addcontentsline{toc}{chapter}{Seznam zkratek}
\begin{description}
  \item[API] Application Programming Interface -- rozhraní pro programování aplikací
  \item[GPIO] General Purpose Input/Output -- univerzální vstupy a výstupy
  \item[PWM] Pulse Width Modulation -- pulzně šířková modulace
\end{description}
\end{verbatim}}

Výsledek (zde bez nadpisu kapitoly):
\begin{description}
	\item[API] Application Programming Interface -- rozhraní pro programování aplikací
	\item[GPIO] General Purpose Input/Output -- univerzální vstupy a výstupy
	\item[PWM] Pulse Width Modulation -- pulzně šířková modulace
\end{description}

Zkratky řaďte abecedně a při prvním použití v textu je vysvětlete,
např.\ \uv{pulzně šířková modulace (PWM)}.
Automaticky vytvářený seznam zkratek (balíček \texttt{glossaries}) popisuje \texttt{README.md};
prodlužuje ale překlad.
